In this section, we encode (co)inductive types in the theory by quantifying over sizes, taking inspiration from the approach of previous work \cite{parametric-quantifier, licentiate}.
We encode these types by constructing initial algebras and terminal coalgebras for polynomial endofunctors; these functors denote the shape of a (co)inductive type.
In particular, inductive types can be described as initial algebras for polynomial endofunctors \cite{containers, homotopy-init}, while coinductive types can be described as terminal coalgebras for polynomial functors, as was done in the construction of coinductive types from inductive types \cite{ahrens_et_al:LIPIcs.TLCA.2015.17}.

\subsection{Polynomial functors}
%In this section, we fix a context $\Gamma$.
As we will be especially interested in endofunctors on the `category' \(U\) of small types and functions, we will omit $\El$ for small types in this section. 
We will also be interested in the `category' of size-indexed small types, where the objects are functions $A : \Size \to \U$, and a morphism between $A, B : \Size \to \U$ is a map $f : \forall i \. (A \, i) \to (B \ i)$, which we will write as $f : A \overset{i}{\to} B$. 

\begin{remark*}
Because we do not assume uniqueness of identity proofs (UIP), $U$ and \(\Size\to U\) are not technically categories, but \((\infty,1)\)-categories \cite{van2011types}.
So, when we talk about functors it would also be more proper to talk about \((\infty,1)\)-functors, which satisfy additional coherence conditions.
Although we suspect that the functors we give here satisfy these coherences, we do not require these conditions for our purpose.
Thus, our internal notions corresponds to those of \emph{wild category theory} \cite{guarded-cubical}, although we will stick to the terms category and functor.
For example, an endofunctor on \(U\) (a functor from \(U\) to \(U\)) consists of the following terms:\begin{align*}
    \F : \U \to \U, \quad\text{and}\quad \F_{A,B} : (A \to B) \to (\F A \to \F B)\text{ for all }A,B:\U,
\end{align*} which satisfy \begin{align*} 
    \phi_A : \F_{A, A}(\id_A) = \id_{FA},  \quad \text{and} \quad
    \phi_{f, g} : \F_{A, C}(g \circ f) = \F_{B, C} \, g \circ \F_{A, B} f.
\end{align*} 
\end{remark*}

We briefly recall polynomial endofunctors, and how they relate to inductive types (W-types) and coinductive types (M-types), which are the types of well-founded trees and non-well-founded trees respectively.
Given a small type $A : U$ and a type $B : A \to U$, the polynomial endofunctor $P_{A, B}$ on the category of small types is defined as:
\begin{align*}
    P_{A,B}(X) &\coloneq  \Sigma{a:A} \. (B(a) \to X ), \\
    P_{A, B}(f) &\coloneq  \lambda \langle a, g\rangle \. \langle a, f \circ g \rangle.
\end{align*}
We can encode W-types and M-types as initial algebras and terminal coalgebras of polynomial endofunctors respectively.
These W-types and M-types can in turn be used to encode all inductive and coinductive types.
We give some examples of (infinitely-branching) inductive types encoded as W-types. More on W-types and M-types, their elimination principles, and  their construction can be found in, for example, \cite{containers, ahrens_et_al:LIPIcs.TLCA.2015.17}.

Given $A : U$ and $B : A \to U$, the associated W-type  \(\rm W x : A.\,B(x) \) of well-founded trees has a single constructor rule:
\begin{center}
\begin{bprooftree}
    \AxiomC{$\Gamma \vdash a : A$}
    \AxiomC{$\Gamma \vdash f : B(a) \to\rm W x : A.\,B(x) $}
    \BinaryInfC{$\Gamma \vdash \mathsf{tree}(a, f) : \rm W x : A.\, B(x) $}
\end{bprooftree}
\end{center}
Intuitively, the type $A$ determines the shapes or constructors of nodes in a well-founded tree, and $B$ determines the branching structure, i.e. arity of each constructor. The term $\mathsf{tree}(a, f)$ can be thought of as the tree with the root node labelled by $a$, and children given by $f$. For example, because the type of natural numbers has two constructors, zero and successor, it is encoded as a W-type by taking the shape to be the boolean type. The zero constructor takes no arguments and the successor constructor takes a single argument. Thus, the type $B : \mathsf{Bool} \to U$ determining the arity should defined by \( B(\mathsf{ff}) = \bot \) and \( B(\mathsf{tt}) = \top \). 

More interestingly, infinitely-branching inductive types can also be encoded as W-types by letting $B$ be infinite. One example is the inductive type $\mathsf{Brw}$ of Brouwer ordinals, which extends the constructors for the natural numbers with a limit constructor $\mathsf{lim} :(\mathbb{N} \to \mathsf{Brw}) \to \mathsf{Brw}$. Informally, assuming an existing type of natural numbers, this is encoded as the type $\rm Wx:A.\,B$ where $A\coloneqq\top+\top+\top$ is a type with three terms, which we may suggestively call $\mathsf{zero}, \mathsf{succ}$, and $\mathsf{lim}$.
We can then define $B : A \to U$ by 
\begin{align*}
    B(\mathsf{zero}) = \bot, \quad B(\mathsf{succ}) = \top, \quad \text{and} \quad B(\mathsf{lim}) = \mathbb{N}.
\end{align*}

Infinitely-branching well-founded trees are also essential in the model of constructive set theory given by Aczel \cite{ACZEL197855}. In this model, the type of all sets is given by the W-type \(W_{A : U} A\), where $U$ is the type of the constructors, and the arity is given by the identity on $U$. The idea for this type is that each set can be viewed as a well-founded tree, where its direct successors are the element of the set. 

\subsection{Size-indexed functors}
In the next sections, we will develop our theory for general functors on small types and size-indexed small types, and show how polynomial functors fit in this framework.
We will call endofunctors on the category of size-indexed small types \emph{size-indexed endofunctors}.
Following existing notation \cite{parametric-quantifier, licentiate}, we define size-indexed endofunctors $\Diamond$ and $\Box$, corresponding to bounded existential and universal quantification.
\begin{proposition}
    We have an endofunctor of the category \(\Size\to U\): 
    \begin{align*}
        &\Diamond : (\Size \to \U) \to (\Size \to \U), &
        &\Diamond_{A,B}:(A \overset{i}{\to} B) \to (\Diamond A \overset{i}{\to} \Diamond B), \\
        &\Diamond \, A  \, i \coloneq   \exists j < i\. A \, j; &
        &\Diamond_{A, B} \, f \, i \, \langle j, a \rangle \coloneq  \brac{j, f \, j \, a}.
    \end{align*}
\end{proposition}

\begin{proposition}
    We have an endofunctor of the category \(\Size\to U\):  
    \begin{align*}
        &\Box : (\Size \to \U) \to (\Size \to \U), &
        &\Box_{A,B} : (A \overset{i}{\to} B) \to (\Box A \overset{i}{\to} \Box B), \\
        &\Box \, A  \, i \coloneq  \forall j < i \. A \, j; &
        &\Box_{A , B} \, f \, i \, g \coloneq  \lambda j \. f \, j \, (g \, j).
    \end{align*}
\end{proposition}
We will write $\Diamond_i \, A$ for $\Diamond \, A \, i \coloneq  \exists j < i \. A \, j$ and similarly $\Box_i \, A$ for $\Box \,  A \, i \coloneq  \forall j < i\. A \, j$.

Given a functor $F$ on small types, we can make use of $\Diamond$ to `lift' $F$ to the category of size-indexed types by precomposing with $\Diamond$:
\begin{align*}
&\F[\Diamond -]  : (\Size \to \U) \to (\Size \to \U), &
&\F[\Diamond -]_{A,B} : (A \overset{i}{\to} B) \to ( \F[\Diamond \, A] \overset{i}{\to} \F[\Diamond \, B]), \\
&\F[\Diamond \, A] \, i \coloneq  \F( \Diamond_i \, A ) \coloneq  F (\exists j < i \. A \, j); &
&\F[\Diamond \, f]_{A, B} \, i \coloneq  F_{\Diamond_i  A, \Diamond_i  B} (\Diamond_{A, B} \, f \, i).
\intertext{Given a functor $F$ on the category of small types describing the shape of an inductive type, the size-indexed functor $\F[\Diamond -]$ should describe the shape of the size-indexed variant of the type. Analogously, \( \F[\Box -] \) is defined by precomposition with \( \Box \):}
    &\F[\Box -]  : (\Size \to \U) \to (\Size \to \U), &
    &\F[\Box -]_{A,B} : (A \overset{i}{\to} B) \to ( \F[\Box \, A] \overset{i}{\to} \F[\Box \, B]), \\
    &\F[\Box \, A] \, i \coloneq  \F( \Box_i \, A ) \coloneq  F (\forall j < i \. A \, j); & 
    &\F[\Box \, f]_{A, B} \, i \coloneq  F_{\Box_i  A, \Box_i  B}(\Box_{A, B} \, f \, i).
\end{align*}

\subsection{Inductive types}
Since we encode inductive type as initial algebras, we recall the notion of an algebra for a endofunctor on $U$ and define it for size-indexed functors.

\begin{definition}[Algebra]
    Given an endofunctor \( F : U \to U \), an \emph{F-algebra} is a pair \( (A, k_A) \) consisting of a type \( A : U \) and a structure map \( k_A : FA \to A \). A map of coalgebras \( (h, s_h) : (A, k_A) \to (B, k_B) \) is a pair of a map \( h : A \to B \) and a path \( s_h : h \circ k_A = k_B \circ Fh \). We write $\FAlg$ for the type of $F$-algebras.
An \emph{initial $F$-algebra} is an algebra \( (\mu F, \In) \), such that for any algebra \( (A, k_A) \), the following type is contractible: 
\[\sum_{h : \mu F \to A}(h \circ \In = k_A \circ Fh).\]
\end{definition}
We define \emph{size-indexed \( F \)-algebras} as algebras for \( F[\Diamond -] \).
Explicitly, a size-indexed \( F \)-algebra is a pair \( (A, k_A) \) where \( A : \Size \to U \) and \( k_A : F[\Diamond A] \sto A \).
% We will write \( F[\Diamond{-}] \)-algebra for \( F[\Diamond -] \)-algebra where convenient, and write $\sFAlg$ for the type of $F[\Diamond{-}]$-algebras. 

To encode size-indexed inductive types, we need to construct initial algebras for size-indexed functors. This is done in the following lemma, where we make use of the fixpoint operator:
\begin{lemma}\label{lemma:inductive-sized-encoding}
    Let \( F : U \to U \) be a functor. The size-indexed functor \( F[\Diamond -] \) has an initial algebra \( (\mu^\Diamond F, \In^\Diamond) \), where the first component is given by:
\[
    \mu^\Diamond F \, i \coloneqq  \fix (\lambda i. \lambda  X. F(\Diamond_i X)) \, i.
\] 
\end{lemma}
\begin{proof}[Proof Sketch]
    We note that the propositional \( \beta \)-rule for the fixpoint operator gives us the following equality:
    \begin{align*}
        \mu^\Diamond F \, i &\coloneqq\fix (\lambda i \. \lambda X\. F(\Diamond_i X)) \, i \\
        &= (\lambda i \. \lambda X\. F(\Diamond_i X)) \, i \, (\lambda r < s \. \fix (\dots) \, r) \\
        &\equiv F(\Diamond_i \, \fix (\lambda s \. \lambda X\. F(\Diamond_s X)) ) \\
        & =: F [\Diamond (\mu^\Diamond F)] \, i
    .\end{align*}
    The structure map \( \In^\Diamond : F[\Diamond (\mu^\Diamond F)] \sto \mu^\Diamond F \) is then given by the pointwise transport operation. Given an algebra \( (A, k_A) \), the term \( \fold^\Diamond \, A \, k_A : \mu^\Diamond F \sto A\) is also defined using the fixpoint operator, and can be shown to be an algebra morphism. Initiality then follows from the contractibility of \( \fix \!\) described in \autoref{lemma:unique-fix}, together with function extensionality.
\end{proof}
Thus, we have a way of constructing size-indexed types, which can be thought of as the ordinal approximations for our inductive types. But our main interest is in showing that $F$ itself has an initial algebra that can be constructed via existential quantification over the initial \( F[\Diamond{-}] \)-algebra.
To this end, we first note the following properties:
\begin{lemma}
    Every \( F \)-algebra \( (A, k_A) \) induces an \( F[\Diamond{-}] \)-algebra \( (TA, k_{TA}) \), where the first component is 
    \(
        TA \coloneqq (\lambda i. A) : \Size \to U.
    \) 
This extends to a functor \( T : \FAlg \to \sFAlg \). 
\end{lemma}
\begin{proof}[Proof Sketch]
    We require a structure map \( k_{TA} : F[\Diamond TA] \sto TA \). Noting that \( i \notin \text{fv}(A) \), we can define the following auxiliary term:
    \begin{align*}
        & \mathsf{extract} : \Pi A:U \. \Diamond T A \sto TA, \\
        & \mathsf{extract}_A \, i \, \langle j, a \rangle \coloneqq  a
    .\end{align*}
    This is well-defined because $TA \, j \equiv A$. As the name suggests, it allows us to extract a term away from the irrelevant size-index.
    This allows us to define \( k_{TA} \, i \coloneqq k_A \circ F(\mathsf{extract}_A \, i)\), since $F(\mathsf{extract}_A \, i) : F[\Diamond_i TA] \to FA$.
    
    Given a algebra map \( f : (A, k_A) \to (B, k_B) \), we lift this to a map between size-indexed algebras \( T\,f : (TA, k_{TA}) \to (TB, k_{TB}) \) defined by \( Tf \coloneqq \lambda i . f \). The fact that this forms an algebra map follows from the fact that $\mathsf{extract}_A$ is natural in $A$.
\end{proof}

We can also define a functor from size-indexed algebras to algebras by existential quantification on the carrier of the size-indexed algebra.
However, this construction only works on the class of functors satisfying the following condition. As we will see, the equivalences we axiomatised ensure that this includes all polynomial functors.

\begin{definition}\label{def:weakly-commute-exist}
    Let $\F$ be an endofunctor on the category of small types. We say that $\F$ \emph{weakly commutes} with the existential type if the following canonical map is an equivalence for any $A : \Size \to \U$:
    \begin{align*}
    \mathsf{can}_A &: \exists i \. F [\Diamond A] \, i \to \F (\exists i .A \, i), & \varphi_A& : \forall i \. (\exists j < i . A \, j \to \exists j . A \, j), \\
    \mathsf{can}_A & \, \langle i, a' \rangle \coloneq  F \, (\varphi_A \, i) \, a' & \varphi_A& \, i \, \brac{j, a} \coloneq  \brac{j, a}.
    \end{align*}
\end{definition}

For functors that satisfy this property, we can define a functorial map $\exists : F[\Diamond{-}]\text-\sf{Alg} \to F\text-\sf{Alg}$ by lifting the existential quantifier to work on algebras.
\begin{proposition}
    Given a functor $F$ on the category of small types that weakly commutes with the existential quantifier, a size-indexed $F$-algebra $(A, k_A)$ gives rise to an $F$-algebra $(\exists i . A \, i, k_{\exists A})$. Moreover, this operation is functorial.
\end{proposition}
\begin{proof}[Proof Sketch]
    We first note that for any map $f : A \sto B$, we can define a map between existentially quantified types as follows:
    \begin{align*}
        \exists f : \exists i . A \, i \to \exists i . B \, i, \\
        \exists f \, \langle i, a \rangle \coloneq  \langle i, f \, i \, a \rangle.
    \end{align*}

    Let $\mathsf{can}^{-1}_A : F(\exists i . A \, i) \to (\exists i . F [\Diamond A] \, i) $ be the equivalence obtained from $F$ weakly commuting with the existential type. We define the map $k_{\exists A} : F (\exists i . A \, i) \to \exists i. A \, i$ by
    \begin{align*}
    k_{\exists A}  \coloneq  \exists \, k_A \circ \mathsf{can}^{-1}_A.
    \end{align*}
    Given an algebra map $f :(A, k_A) \to (B, k_B)$, we define an algebra map $(\exists i . A \, i, k_{\exists A}) \to (\exists i . B \, i, k_{\exists B})$ as $\exists f : \exists i . A \, i \to \exists i . B \, i$.
    The fact that this constitutes an algebra morphism follows from $\mathsf{can}_A$ being natural in $A$. Thus, we have a functor $\exists : \sFAlg \to \FAlg$.
\end{proof}

Finally, we are able to prove our desired result, namely that for functors that weakly commute with the existential type, their initial algebra is obtained by existential quantification over the initial size-indexed algebra.
\begin{theorem}\label{theorem:inductive-construction}
    Let $\F$ be a functor on the category of small types that weakly commutes with the existential type. Then $\F$ has an initial algebra with first component:
    \[
    \mu  \F \coloneq  \exists i . (\mu^\Diamond F) \, i.
    \]
\end{theorem}
\begin{proof}[Proof Sketch]
    The categorical intuition is that this follows from the fact that the functorial map $\exists : \sFAlg \to \FAlg$ of the previous proposition is the left adjoint of $T : \FAlg \to \sFAlg$.
    Since left adjoints preserve initial objects, we find that $\exists i. (\mu^\Diamond F) \, i$ is the initial algebra of $F$, with the structure map $\In \coloneq  \exists \, \In^\Diamond \circ \mathsf{can}^{-1}_{\mu^\Diamond F}$, which is also an equivalence.

    We can work this out in more detail: for any $A : \Size \to U$ and $B : U$, there is an isomorphism between maps of type $A \sto TB$ and those of type $\exists i . A \, i \to B$, given by the currying isomorphism. This extends to an isomorphism between size-indexed algebra maps of type $(A, k_A) \to (TB, k_{TB})$ and algebra maps of type $(\exists i . A \, i, k_{\exists A}) \to (B, k_B)$. Thus, defining an algebra map $(\mu \F, \In) \to (B, k_B)$ is equivalent to defining a size-indexed algebra map $(\mu^\Diamond F, \In^\Diamond) \to (TB, k_{TB})$. This map is uniquely given by the initiality of $(\mu^\Diamond F, \In^\Diamond)$.  
\end{proof}

Lastly, we show that all polynomial endofunctors on small types satisfy the condition of weakly commuting with the existential quantifier.
The equivalences that we have axiomatised entail the following chain of equivalences, which corresponds precisely to the canonical map denoted in \autoref{def:weakly-commute-exist}, noting that $i \notin \text{fv}(A)$ and $i \notin \text{fv}(B(a))$ for $a : A$:
\begin{align*}
    \exists i . P_{A,B}(\Diamond X) \, i &\coloneq  \exists i \. \Sigma {a : A}  \. (B(a) \to \exists j < i . X \, j) \\
    &\simeq \Sigma {a : A} \. (\exists i . B(a) \to \exists j < i . X \, j)\\
    &\simeq \Sigma {a : A} \. (B(a) \to \exists i . \exists j < i . X \, j)\\
    &\simeq \Sigma {a : A} \. (B(a) \to \exists i . X \, i) =: P_{A,B}(\exists i . X \, i).
\end{align*}
Since all polynomial functors (both finitely- and infinitely-branching) weakly commute with the existential quantifier, their initial algebra is constructed via existential quantification over the corresponding initial size-indexed algebra. This allows us to construct W-types via quantification over size-indexed types.

\subsection{Coinductive types}
The coalgebra construction is similar to that of algebras. We define size-indexed coalgebras and construct their terminal coalgebras using the fixpoint operator. This allows us to obtain terminal coalgebras for functors satisfying an analogous property: weakly commuting with the universal quantifier. This is done by showing that the universal quantifier lifts to a functor from size-indexed coalgebras to coalgebras, which is right adjoint to the inclusion map from coalgebras to size-indexed coalgebras. 
The equivalences allow us to show that all polynomial functors weakly commute with the universal quantifier.
The details are given in \autoref{appendix:encoding}.

